\chapter{Addressing Modes}
\label{ch:addressing-modes}

\section{Overview}

The SIRC-1 CPU supports seven addressing modes that specify how operands are accessed. These modes determine whether an
operand is an immediate value, register contents, an address-register pair named directly, or data in memory at a
computed address. The 8-bit short-immediate encoding used by ALU shift-and-add forms is an encoding-width variant of
Immediate addressing, not a separate mode; Table~\ref{tab:addr-mode-forms} maps every operand form the assembler
accepts onto one of the seven modes.

\begin{table}[H]
	\centering
	\begin{tabular}{lp{8cm}}
		\toprule
		\textbf{Mode}           & \textbf{Description}                                                 \\
		\midrule
		Immediate               & Operand is a constant value encoded in the instruction               \\
		Register Direct         & Operand is the contents of a specified register                      \\
		Address Register Direct & Operand names an address-register pair directly, with no indirection \\
		Indirect Immediate      & Memory address = address register + immediate offset                 \\
		Indirect Register       & Memory address = address register + register offset                  \\
		Post-Increment          & Indirect register, then increment address register                   \\
		Pre-Decrement           & Decrement address register, then use as indirect                     \\
		\bottomrule
	\end{tabular}
	\caption{Addressing Modes Summary}
	\label{tab:addr-modes}
\end{table}

\begin{table}[H]
	\centering
	\scriptsize
	\begin{tabularx}{\textwidth}{>{\raggedright\arraybackslash}p{5.5cm}>{\raggedright\arraybackslash}p{3.2cm}>{\raggedright\arraybackslash}X}
		\toprule
		\textbf{Assembler operand form}                 & \textbf{Example}                  & \textbf{Addressing mode}                                                                             \\
		\midrule
		Immediate                                       & \texttt{\#100}                    & Immediate                                                                                            \\
		Immediate, 8-bit shift-qualified                & \texttt{\#10, LSL \#2}            & Immediate (encoding-width variant; see Short Immediate Addressing)                                   \\
		Register Direct                                 & \texttt{r1}                       & Register Direct                                                                                      \\
		Address Register Direct                         & \texttt{a}                        & Address Register Direct                                                                              \\
		Indirect Register Displacement                  & \texttt{(r2, a)}                  & Indirect Register                                                                                    \\
		Indirect Immediate Displacement                 & \texttt{(\#8, a)}, \texttt{(a)}   & Indirect Immediate                                                                                   \\
		Indirect Register Displacement, Post-Increment  & \texttt{(r2, s)+}                 & Post-Increment                                                                                       \\
		Indirect Immediate Displacement, Post-Increment & \texttt{(\#2, s)+}, \texttt{(s)+} & Post-Increment                                                                                       \\
		Indirect Register Displacement, Pre-Decrement   & \texttt{-(r2, s)}                 & Pre-Decrement                                                                                        \\
		Indirect Immediate Displacement, Pre-Decrement  & \texttt{-(\#2, s)}, \texttt{-(s)} & Pre-Decrement                                                                                        \\
		Shift Definition                                & \texttt{LSL \#2}                  & Not an addressing mode; modifies the register operand it follows (Chapter~\ref{ch:shift-operations}) \\
		\bottomrule
	\end{tabularx}
	\caption{Assembler Operand Forms Mapped to Addressing Modes}
	\label{tab:addr-mode-forms}
\end{table}

\section{Addressing Mode Matrix}

\begin{table}[H]
	\centering
	\scriptsize
	\begin{tabularx}{\textwidth}{>{\raggedright\arraybackslash}p{1.8cm}>{\raggedright\arraybackslash}X>{\raggedright\arraybackslash}X>{\raggedright\arraybackslash}p{1.9cm}>{\raggedright\arraybackslash}X}
		\toprule
		\textbf{Mode}           & \textbf{Syntax}                                                             & \textbf{Effective value or address}                                                                                  & \textbf{Memory access}                  & \textbf{Legal families}                              \\
		\midrule
		Immediate               & \texttt{\#imm16}, or \texttt{\#imm8, shift} (short-immediate width variant) & 16-bit encoded value; the 8-bit short-immediate form is zero-extended, and its shift applies to the register operand & None                                    & ALU, LOAD-immediate, COPI, branch meta-instructions  \\
		Register Direct         & \texttt{rN}                                                                 & 16-bit register value                                                                                                & None                                    & ALU, LOAD register-copy, COPR, register-offset forms \\
		Address Register Direct & \texttt{addr}                                                               & \texttt{addr.high : addr.low}, used directly                                                                         & None                                    & LJMP/LJSR src, LDEA/LDEL dest                        \\
		Indirect Immediate      & \texttt{(\#offset, addr)}, \texttt{(addr)}                                  & \texttt{addr.high : (addr.low + offset)}                                                                             & LOAD reads; STOR writes; LDEA/LDEL none & Memory, LDEA, LDEL                                   \\
		Indirect Register       & \texttt{(rO, addr)}                                                         & \texttt{addr.high : (addr.low + rO)}                                                                                 & LOAD reads; STOR writes; LDEA/LDEL none & Memory, LDEA, LDEL                                   \\
		Post-Increment          & \texttt{(\#offset, addr)+}, \texttt{(addr)+}, \texttt{(rO, addr)+}          & Address is calculated from original \texttt{addr}; \texttt{addr.low += 1} after access                               & LOAD reads; LDEL none                   & LOAD, LDEL, LJSR (assembles to LDEL)                 \\
		Pre-Decrement           & \texttt{-(\#offset, addr)}, \texttt{-(addr)}, \texttt{-(rO, addr)}          & \texttt{addr.low -= 1}; address then uses decremented \texttt{addr.low + offset}                                     & STOR writes; LDEA none                  & STOR, LDEA (and LDEA-derived meta-instructions)      \\
		\bottomrule
	\end{tabularx}
	\caption{Addressing Mode Side Effects and Legal Instruction Families}
	\label{tab:addressing-mode-matrix}
\end{table}

\section{Legal Modes by Instruction Family}

Table~\ref{tab:addressing-mode-matrix} answers "what does this addressing-mode syntax mean, and which instruction families can use it in general?" Table~\ref{tab:legal-modes-by-family} answers the reverse question, "for this specific instruction, exactly which operand forms are legal?" The two tables must agree; consult the one that matches the question you are asking.

\begin{table}[H]
	\centering
	\scriptsize
	\begin{tabularx}{\textwidth}{>{\raggedright\arraybackslash}p{2.2cm}>{\raggedright\arraybackslash}X>{\raggedright\arraybackslash}X}
		\toprule
		\textbf{Instruction family}    & \textbf{Legal operand forms}                                                                                                                                                                       & \textbf{Explicit exclusions and notes}                                                                                               \\
		\midrule
		ALU result instructions        & \texttt{OPI rD, \#imm16}; \texttt{OPI rD, \#imm8, shift}; \texttt{OPR rD, rS1, rS2[, shift]}; register shorthand \texttt{OPR rD, rS2}.                                                             & No memory operands. Full 16-bit immediate forms do not encode shifts.                                                                \\
		ALU test instructions          & \texttt{CMPI/TSAI/TSXI rS, \#imm16}; \texttt{CMPI/TSAI/TSXI rS, \#imm8, shift}; \texttt{CMPR/TSAR/TSXR rS1, rS2[, shift]}.                                                                         & No destination write-back. No memory operands.                                                                                       \\
		LOAD value                     & \texttt{LOAD rD, \#imm16}; \texttt{LOAD rD, rS}.                                                                                                                                                   & No public short-immediate LOAD syntax. Direct register-copy forms do not accept shift suffixes.                                      \\
		LOAD memory                    & \texttt{LOAD rD, (\#offset, addr)}; \texttt{LOAD rD, (rO, addr)}; \texttt{LOAD rD, (\#offset, addr)+}; \texttt{LOAD rD, (rO, addr)+}.                                                              & No pre-decrement LOAD form. No shift suffix on any \mnemonic{LOAD} form.                                                              \\
		STOR memory                    & \texttt{STOR (\#offset, addr), rS}; \texttt{STOR (rO, addr), rS[, shift]}; \texttt{STOR -(\#offset, addr), rS}; \texttt{STOR -(rO, addr), rS[, shift]}.                                            & No post-increment STOR form. Shift suffix is only available on register-displacement \mnemonic{STOR} forms.                          \\
		LDEA                           & \texttt{LDEA dest, (\#offset, src)}; \texttt{LDEA dest, (rO, src)}; \texttt{LDEA dest, -(\#offset, src)}; \texttt{LDEA dest, -(rO, src)}.                                                          & No post-increment LDEA form. No shift suffixes.                                                                                      \\
		LDEL                           & \texttt{LDEL dest, (\#offset, src)}; \texttt{LDEL dest, (rO, src)}; \texttt{LDEL dest, (\#offset, src)+}; \texttt{LDEL dest, (rO, src)+}.                                                          & No pre-decrement LDEL form. No shift suffixes.                                                                                       \\
		Control-flow meta-instructions & \texttt{BRAN/BRSR \#disp}; \texttt{BRAN/BRSR @label}; \texttt{RETS}; \texttt{LJMP src[, \#offset|rO]}; \texttt{LJSR src[, \#offset|rO]}; \texttt{LJSR (\#offset, src)+}; \texttt{LJSR (rO, src)+}. & \texttt{LJMP (src)} and \texttt{LJSR (src)} are not public syntax; use \texttt{LJMP src} or \texttt{LJSR src} for zero displacement. \\
		Coprocessor                    & \texttt{COPI \#imm16}; \texttt{COPR rS}; standard coprocessor meta-instructions as documented in Chapter~\ref{ch:coprocessor}.                                                                     & No hidden destination register syntax, no memory operands, no shifts, and no status-update override suffixes.                        \\
		\bottomrule
	\end{tabularx}
	\caption{Addressing Mode Legality by Instruction Family}
	\label{tab:legal-modes-by-family}
\end{table}

\section{Zero-Displacement Shorthand}

Immediate-displacement indirect forms may omit the displacement when it is zero. The following forms are exact
assembler-level equivalents:

\begin{table}[H]
	\centering
	\begin{tabular}{lll}
		\toprule
		\textbf{Shorthand} & \textbf{Canonical form} & \textbf{Meaning}                  \\
		\midrule
		\texttt{(addr)}    & \texttt{(\#0, addr)}    & Use the selected address register \\
		\texttt{(addr)+}   & \texttt{(\#0, addr)+}   & Use the address, then increment   \\
		\texttt{-(addr)}   & \texttt{-(\#0, addr)}   & Decrement, then use the address   \\
		\bottomrule
	\end{tabular}
	\caption{Zero-Displacement Shorthand}
	\label{tab:zero-displacement-shorthand}
\end{table}

The shorthand does not change which instruction families are legal. For example, \texttt{LOAD r1, (a)+} is legal
because \mnemonic{LOAD} has a post-increment immediate-displacement form, but \texttt{STOR (a)+, r1} is not legal
because \mnemonic{STOR} has no post-increment form. Likewise, \texttt{STOR -(s), r1} is legal because \mnemonic{STOR}
has a pre-decrement form, but the corresponding pre-decrement \mnemonic{LOAD} form is not legal.

Register-displacement forms do not have a corresponding shorthand: \texttt{(rO, addr)}, \texttt{(rO, addr)+}, and
\texttt{-(rO, addr)} always name the displacement register explicitly.

\section{Operand Categories}

Instruction descriptions use these operand categories:

\begin{description}
	\item[Register value] A 16-bit general-purpose, status, or address-component register operand used directly by an ALU, load,
	      store, or coprocessor instruction.
	\item[Immediate value] A constant encoded in the instruction. Full immediates are 16 bits; short immediates are 8 bits.
	\item[Memory source] An effective address read by \mnemonic{LOAD}.
	\item[Memory destination] An effective address written by \mnemonic{STOR}.
	\item[Branch target] An effective address written to the program-counter pair \reg{p}.
	\item[Address destination] An address-register pair written by \mnemonic{LDEA} or \mnemonic{LDEL}.
	\item[Alterable destination] Any destination register, address-register pair, or memory location that the instruction can
	      modify when its condition is true.
\end{description}

\section{Common Address Calculation Rules}

For indirect, post-increment, pre-decrement, branch, jump, and effective-address forms, address arithmetic is performed
on the low word of the selected address-register pair. The high word supplies the segment and is preserved unless the
instruction explicitly writes an address-register pair destination.

Post-increment and pre-decrement change the selected address register by exactly one word. The displacement or index
register participates only in the effective address calculation; it does not change the auto-update amount.

When \reg{p} is used as the source address register, the value used for address calculation is the address of the
instruction currently executing. PC-relative displacements are therefore relative to the address of the branch
instruction itself, not to the next instruction; a displacement of 2 targets the following instruction and a
displacement of 0 re-executes the branch. \mnemonic{BRSR} and \mnemonic{LDEL} still link to \reg{pl} + 2.

\section{Per-Mode Encoding Reference}

Every SIRC-1 instruction is a fixed 32-bit (2-word) quantity regardless of addressing mode -- unlike architectures
with variable-length extension words per mode, there is no mode whose word count differs from any other's. Each box
below states this explicitly for completeness against that comparison point, alongside the generation formula,
assembler syntax, and field encoding for that mode.

\begin{addrmodebox}{Immediate Addressing}
	\textbf{Generation:} Operand = the 16-bit value encoded directly in the instruction.

	\textbf{Assembler Syntax:}
	\begin{lstlisting}
ADDI r1, #100         ; 16-bit immediate
ADDI r2, #-50         ; Signed 16-bit immediate
	\end{lstlisting}

	\textbf{Field Encoding:} Immediate format; Immediate field, bits 21--6 (16 bits, arithmetic instructions
	interpret this field according to Chapter~\ref{ch:data-representation}); Register field, bits 25--22
	(destination).

	\textbf{Instruction Word Count:} 2 words (fixed).

	\textbf{Usage:} Loading constant values; arithmetic with known values; bit manipulation with masks.
\end{addrmodebox}

\begin{addrmodebox}{Register Direct Addressing}
	\textbf{Generation:} Operand = the 16-bit value currently stored in the specified register. This is the most
	common mode for ALU operations in a load/store architecture.

	\textbf{Assembler Syntax:}
	\begin{lstlisting}
ADDR r1, r2, r3   ; r1 = r2 + r3
ADDR r4, r5       ; r4 = r4 + r5
	\end{lstlisting}

	\textbf{Field Encoding:} Register format; Register 1, bits 25--22 (destination); Register 2, bits 21--18
	(source A); Register 3, bits 17--14 (source B).

	\textbf{Instruction Word Count:} 2 words (fixed).

	\textbf{Usage:} All ALU operations between registers; register-to-register data movement; testing register
	values.
\end{addrmodebox}

\begin{addrmodebox}{Address Register Direct Addressing}
	\textbf{Generation:} Operand = the full 32-bit contents (high word and low word) of the named address-register
	pair (\reg{l}, \reg{a}, \reg{s}, or \reg{p}), with no indirection. \mnemonic{LDEA} and \mnemonic{LDEL} use this
	form to select the destination pair that receives the computed address; \mnemonic{LJMP} and \mnemonic{LJSR}
	use it to select the source pair whose contents become the new value of \reg{p}.

	\textbf{Assembler Syntax:}
	\begin{lstlisting}
LDEA a, (#4, s)       ; a = s + 4 (destination named directly)
LJMP a                ; p = a (source named directly)
	\end{lstlisting}

	\textbf{Field Encoding:} For \mnemonic{LDEA}/\mnemonic{LDEL} (opcodes \opcode{18}--\opcode{1F}), a 2-bit
	address-register-pair index in bits 23--22 (00 = l, 01 = a, 10 = s, 11 = p); bits 25--24 are unused for this
	purpose. \mnemonic{LJMP}/\mnemonic{LJSR} are meta-instructions that lower to this same encoding with \reg{p}
	as the named pair.

	\textbf{Instruction Word Count:} 2 words (fixed).

	\textbf{Usage:} Selecting which address-register pair an effective-address instruction writes; jumping or
	calling through an address already held in a register pair.
\end{addrmodebox}

\begin{addrmodebox}{Indirect Immediate Addressing}
	\textbf{Generation:}
	\begin{equation*}
		\text{EA} = \text{AddressRegister.high} : (\text{AddressRegister.low} + \text{Offset})
	\end{equation*}
	The address register can be \reg{l} (link), \reg{a} (address), \reg{s} (stack), or \reg{p} (program counter).
	Offset is a signed two's-complement value. When \texttt{SR.A} (Trap on Address Overflow) is set, the Segment
	Overflow Fault (Chapter~\ref{ch:exceptions}) fires only if the low-word addition wraps; a negative offset that
	keeps the sum within the low word's range does not fault.

	\textbf{Assembler Syntax:}
	\begin{lstlisting}
LOAD r1, (#8, a)      ; Load from address (a + 8)
STOR (#-4, s), r2     ; Store to address (s - 4)
LDEA p, (#100, p)     ; Jump relative (pc = pc + 100)
	\end{lstlisting}

	\textbf{Field Encoding:} Memory/control format, an even opcode in \opcode{10}--\opcode{1F}; 2-bit
	address-register-pair index in bits 23--22; 16-bit signed immediate offset in bits 21--6.

	\textbf{Instruction Word Count:} 2 words (fixed).

	\textbf{Usage:} Accessing structure fields (fixed offset from base); array element access (constant index);
	stack frame access (local variables); position-independent code (PC-relative addressing).
\end{addrmodebox}

\begin{addrmodebox}{Indirect Register Addressing}
	\textbf{Generation:}
	\begin{equation*}
		\text{EA} = \text{AddressRegister.high} : (\text{AddressRegister.low} + \text{IndexRegister})
	\end{equation*}
	Identical to Indirect Immediate Addressing except the offset is taken from a register at run time instead of
	being encoded as an immediate value.

	\textbf{Assembler Syntax:}
	\begin{lstlisting}
LOAD r1, (r2, a)      ; Load from address (a + r2)
STOR (r4, s), r3      ; Store to address (s + r4)
	\end{lstlisting}

	\textbf{Field Encoding:} Memory/control format, an odd opcode in \opcode{10}--\opcode{1F}; 2-bit
	address-register-pair index in bits 23--22; index register in bits 17--14. Bits 21--18 are unused for
	\mnemonic{LOAD} and hold the (shifted) source value for \mnemonic{STOR}.

	\textbf{Instruction Word Count:} 2 words (fixed).

	\textbf{Usage:} Array indexing with a variable index; table lookups; dynamic offset calculations.
\end{addrmodebox}

\begin{addrmodebox}{Post-Increment Addressing}
	\textbf{Generation:}
	\begin{enumerate}
		\item Compute EA = AddressRegister.high : (AddressRegister.low + Offset)
		\item Perform the memory operation (load)
		\item AddressRegister.low = AddressRegister.low + 1
	\end{enumerate}
	Particularly useful for iterating through arrays or implementing stack operations.

	\textbf{Assembler Syntax:}
	\begin{lstlisting}
LOAD r1, (r2, s)+     ; Load from (s + r2), then s = s + 1
LOAD r1, (#2, s)+     ; Load from (s + 2), then s = s + 1
	\end{lstlisting}

	\textbf{Field Encoding:} Selected by the opcode itself, not a separate operand field: opcode bit 1 = 1
	selects auto-update, and bit 0 selects an immediate (\opcode{16}) or register (\opcode{17}) offset.

	\textbf{Instruction Word Count:} 2 words (fixed).

	\textbf{Usage:} Stack pop operations; forward array traversal; buffer reading.
\end{addrmodebox}

\begin{addrmodebox}{Pre-Decrement Addressing}
	\textbf{Generation:}
	\begin{enumerate}
		\item AddressRegister.low = AddressRegister.low - 1
		\item Compute EA = AddressRegister.high : (AddressRegister.low + Offset)
		\item Perform the memory operation (store)
	\end{enumerate}
	The complement of post-increment; useful for stack push operations.

	\textbf{Assembler Syntax:}
	\begin{lstlisting}
STOR -(r2, s), r1     ; s = s - 1, then store to (s + r2)
STOR -(#2, s), r3     ; s = s - 1, then store to (s + 2)
	\end{lstlisting}

	\textbf{Field Encoding:} Selected by the opcode itself, not a separate operand field: opcode bit 1 = 1
	selects auto-update, and bit 0 selects an immediate (\opcode{12}) or register (\opcode{13}) offset.

	\textbf{Instruction Word Count:} 2 words (fixed).

	\textbf{Usage:} Stack push operations; backward array traversal; buffer writing.
\end{addrmodebox}

\begin{addrmodebox}{Short Immediate Addressing}
	\textbf{Generation:} Operand = an 8-bit value encoded in the instruction, zero-extended to 16 bits before ALU
	arithmetic. This is a width variant of Immediate addressing, not a separate addressing mode; see
	Table~\ref{tab:addr-mode-forms}. The shift applies to the register operand, not to the immediate constant.

	\textbf{Assembler Syntax:}
	\begin{lstlisting}
ADDI r1, #10, LSL #2  ; r1 = (r1 << 2) + 10
SUBI r2, #1, LSL #8   ; r2 = (r2 << 8) - 1
	\end{lstlisting}

	\textbf{Field Encoding:} Short-Immediate-with-Shift format; Immediate field, bits 21--14 (8 bits); Register
	field, bits 25--22 (destination).

	\textbf{Instruction Word Count:} 2 words (fixed).

	\textbf{Usage:} Loading constant values; arithmetic with known values; bit manipulation with masks.
\end{addrmodebox}

\section{Address Register Selection}

Many addressing modes require specifying which address register pair to use. The 2-bit address register field encodes:

\begin{table}[H]
	\centering
	\begin{tabular}{ccl}
		\toprule
		\textbf{Bits} & \textbf{Symbol} & \textbf{Register Pair}    \\
		\midrule
		00            & l               & Link Register (lh, ll)    \\
		01            & a               & Address Register (ah, al) \\
		10            & s               & Stack Pointer (sh, sl)    \\
		11            & p               & Program Counter (ph, pl)  \\
		\bottomrule
	\end{tabular}
	\caption{Address Register Encoding}
	\label{tab:addr-reg-encoding}
\end{table}

\section{Addressing Mode Examples}

\subsection{Example 8-1: Structure Access}

\begin{lstlisting}
; Assume 'a' points to a struct:
; struct Point {
;   word 0: int16_t x
;   word 1: int16_t y
; }

LOAD r1, (#0, a)      ; r1 = point.x
LOAD r2, (#1, a)      ; r2 = point.y
\end{lstlisting}

\subsection{Example 8-2: Array Iteration}

\begin{lstlisting}
; Iterate array of 16-bit values
; Assume 'a' points to array start
; r7 = counter
:loop
  LOAD r1, (#0, a)+   ; Load element, advance by 1
  ; ... process r1 ...
  SUBI r7, #1
  BRAN|!= @loop
\end{lstlisting}

\subsection{Example 8-3: Stack Frame}

\begin{lstlisting}
:begin_subroutine

; Function prologue - save registers
STOR -(#0, s), r1
STOR -(#0, s), r2
STOR -(#0, s), r3

; Access parameters (above saved regs)
LOAD r4, (#3, s)      ; First parameter
LOAD r5, (#4, s)      ; Second parameter

; Function epilogue - restore registers
LOAD r3, (#0, s)+
LOAD r2, (#0, s)+
LOAD r1, (#0, s)+
RETS
\end{lstlisting}

\section{Addressing Mode Restrictions}

\subsection{Privilege Mode Restrictions}

Protected mode restricts writes to the high word of an address register:

\begin{itemize}
	\item \textbf{Direct high-register writes}: Trigger a privilege violation fault
	\item \textbf{Address-register-pair write-back}: Allowed only when the high word would remain unchanged; otherwise triggers a privilege violation fault
	\item \textbf{Post-increment/Pre-decrement}: Always allowed in protected mode; the auto-update writes only the low word of the pair and never carries into the high word
\end{itemize}

\subsection{Alignment}

\begin{itemize}
	\item All memory addresses refer to 16-bit words
	\item There are no alignment restrictions for general memory I/O
	\item All instructions span 2 words and must begin at an even word address
	\item Unaligned instruction fetches (odd word address) trigger an alignment fault exception
\end{itemize}
